\section{Sufficient global bounds and a concrete research agenda}
\label{sec:global-agenda}

The local theorems expose several quantities that could support a global
bound. At present these quantities describe an execution; they are not
proved restrictions on every diagram. The next task is to convert selected
observations into uniform bounds, while retaining an algorithm that always
finishes with a correct decision.

\subsection{A sufficient accounting principle}

Let $S_*(n)$ bound the complete bit length of every state, query input,
output, and retained certificate record in an execution. Let $Q_*(n)$ bound
the number of charged operations, including failed branches, hierarchy
repairs, canonicalization, certificate verification and changes of
representation. Fixed constants in local polynomial algorithms must be
uniform over the allowed input family. In particular, a handle-type catalog
or a frontier alphabet whose size grows with $n$ cannot be hidden inside
such a constant.

\begin{proposition}[Conditional composition bound]
\label{prop:global-accounting}
Suppose a sound and complete recognition algorithm has preprocessing time
$P(n)$, makes at most $Q_*(n)$ calls, and each call costs at most
$C(1+S_*(n))^a$ bit operations for fixed $C,a$. If, for fixed
$\alpha,\beta,\gamma\ge1$,
\[
\begin{split}
 Q_*(n)&\le\exp(O(\log^{\alpha}(n+2))),\\
 S_*(n)&\le\exp(O(\log^{\beta}(n+2))),\\
 P(n)&\le\exp(O(\log^{\gamma}(n+2))),
\end{split}
\]
then its total time is
$\exp(O(\log^{\max(\alpha,\beta,\gamma)}(n+2)))$.
In particular, exponents at most two suffice for $n^{O(\log n)}$.
\end{proposition}

\begin{proof}
The total time is at most $P(n)+C Q_*(n)(1+S_*(n))^a$.
Taking logarithms of the second term adds the bounds for $\log Q_*$
and $a\log(1+S_*)$. A fixed finite sum of these logarithmic powers is
bounded by a constant times the largest one. Adding preprocessing does
not change that largest exponent.
\end{proof}

The proposition is an accounting criterion, not a theorem that its
hypotheses hold for the present recognizer. Polynomial time for a local
query does not bound $Q_*$ or $S_*$. An inconclusive budget exit does not
supply completeness. Nor does a short successful certificate show that a
deterministic search finds it within the same bound. All expanded data
that must actually be written count towards the relevant size and time.

\subsection{Hierarchy depth, restarts and branching are separate parameters}

Write $H$ for maximum hierarchy length and $G$ for maximum pattern
complexity. In this notation Lackenby's Proposition~9.1 gives the iteration
bound $H(G+1)^H$ \cite{lackenby2026}. Hence a sufficient restriction for
this part of the calculation is
\[
                 H\log(G+1)+\log(H+1)=O(\log^2 n).
\]
For example, $H=O(\log n)$ and $G=n^{O(1)}$ suffice. If instead
$H=O(\log^2 n)$ and $G=n^{O(1)}$, this particular bound gives
$\exp(O(\log^3 n))$: quasi-polynomial in the broader sense, but not the
specific $n^{O(\log n)}$ target. A polynomial bound on the bit length of
$G$ alone does not imply $G=n^{O(1)}$.

The refined construction in Section~10 controls hierarchy length linearly
in the initial $q$-complexity. It does not by itself give logarithmic
depth or a quasi-polynomial bound on the restart expression. Conversely,
replacing a costly surface-topology query by AHT lowers the cost of one
step without lowering the number of lexicographic repairs. These are
different improvements and require different proofs.

If a proposed construction has at most $B(n)$ possible children at depth
at most $H(n)$, its search tree can contain as many as
$\sum_{j=0}^{H(n)}B(n)^j$ nodes. Polynomial branching and logarithmic
depth are sufficient for a quasi-polynomial tree-size bound. Polynomial
depth is not. Revisiting or modifying a node requires an additional
amortized bound, even if the number of distinct nodes is controlled.
Memoization helps only after proving that the canonical state retains
every relevant boundary marking and that the number of such states is
small. A state with polynomially many bits can still have exponentially
many possible values.

\paragraph{Falsifiable hierarchy question.}
For the multisurface and generalized-Heegaard construction, identify a
potential that pays both for ordinary progress and for discarding later
surfaces after a repair. Can every long execution produce a certified
Cheeger-region simplification before the restart cost exceeds a
polylogarithmic exponent? Useful negative instances would keep the
recorded local complexities small while forcing many successive repairs.
Such examples would refute a proposed potential, even if all topology
queries themselves were fast. The announced acceleration framework
\cite{lackenbytalk} motivates this question; the present kernels do not
implement its global argument.

\subsection{From normal-coordinate topology to certified cutting}

The current normal adapter supplies connectivity and total topological
statistics for a supplied surface. A compressed cutting primitive has a
larger contract. It must return the complementary handle structure,
the parallelity components, the topology and twisting of their base
surfaces, and the attachment of their vertical boundaries. These are the
kinds of output required by Lackenby's Theorem~14.4
\cite{lackenby2026}. Matching component counts do not determine these
attachments. Even matching unmarked surface types does not determine a
map relative to a marked boundary.

The first implementation target is weighted AHT with grouped multiplicity
records, followed by a compiler from those records to the required cut
incidence data. The weighted orbit theorem is classical \cite{aht}; its
implementation would extend the current module rather than establish a
new existence theorem. Its output must remain grouped when a binary vector
describes exponentially many equal components. Explicitly listing one
record per component defeats the compression.

A useful contract for testing is a commuting reconstruction diagram:
on small inputs, expand the original surface and perform an independent
ordinary cut; expand the compressed output; compare the resulting marked
cell structures. Include twisted bundles, annular interfaces,
self-identifications, boundary-parallel discs and disconnected surfaces.
For large inputs, retain locally checkable construction records instead
of treating a numerical topology summary as a reconstruction proof.

An equally important target is a verified diagram-to-exterior adapter.
It must specify how the peripheral torus corresponds to the knot and how
that correspondence changes during simplification. A normal disc whose
boundary is essential in a supplied torus certifies compression of that
torus. Turning it into a certificate about the input PD diagram requires
the exterior correspondence. A complete global theorem also needs a
method of finding suitable surfaces within its work budget.

\subsection{Frontier and representation bounds}

Let $w$ be the number of interface positions in a scanner state, $\tau$
the number of retained scalar types, $\mu=\max_v s_v$, $E$ the number of
nonzero typed blocks, $\rho=\max_{u,v}\rho_{uv}$, and $B$ the largest
packed morphism bit length. The elementary inequalities
\[
              V\le\tau\mu^2,\qquad
              P_{\mathrm{span}}\le E\rho\mu^2
\]
describe how these parameters enter the solver. They do not bound the
parameters in terms of $n$. A small coefficient span does not bound
multiplicity, and small frontier width does not by itself bound the
size of all coefficients or the number of repeated states.

One sufficient approach is to prove a uniform encoding bound of the form
$S_*\le\exp(O(w\log(w+2)))\poly(n)$ and a suitable bound on all visits,
then establish $w\log(w+2)=O(\log^2 n)$ for the chosen decomposition or
for an exact quotient of its states. This is a proposed route, not a
property established for the present scanner. A square-root-width
separator bound, even if available for a particular diagram graph, would
give a much larger exponent in this model and would not meet the target.

The experiment should record $w,\tau,\mu,E,\rho,B,V,V_F$, reconstruction
cost, and the number of repeated states together. Search for families in
which one proposed controlling parameter stays fixed while another grows.
In particular, a useful challenge to the forest proposal is a family with
small $V_F$ but costly dense reconstruction in the original $V$ variables.
The current bounded-candidate policy requires that reconstruction, so it
cannot be excluded from the accounting.

\subsection{A tractable commutant family and its limitations}

There is a clean special case in which the forest construction supplies
explicit projections, rather than merely a smaller linear system.

\begin{proposition}[Scalar transports give explicit repeated summands]
\label{prop:scalar-transport-splitting}
Consider an entire finite differential component with no incident
differential blocks omitted. Suppose every scalar type has multiplicity
$s$ and the selected invertible forest connects all its types. Use the
transport matrices $P_v$ from Theorem~\ref{thm:forest}. The complete
commutant is the common centralizer in $M_s(\F)$ of every transported
coefficient matrix $P_v^{-1}A_iP_u$.

If every such matrix is a scalar multiple of $I_s$, the component is
isomorphic to a direct sum of $s$ identical multiplicity-one typed
complexes. Explicit compatible matrix units in the original coordinates
are
\[
                   E^{(v)}_{ij}=P_v E_{ij}P_v^{-1}.
\]
In particular, the $s$ diagonal units are orthogonal rank-one projectors
whose sum is the identity. If additional attachment maps are required in
the splitting contract, the assertion applies when they also respect
these summands; this must be verified, not inferred from the differential.
\end{proposition}

\begin{proof}
There is one root matrix $X$ because the forest is connected. Theorem
\ref{thm:forest} changes every original equation into $X\widetilde A_i
=\widetilde A_iX$, proving the centralizer statement. Under the scalar
hypothesis every $X\in M_s(\F)$ satisfies these equations. After changing
the basis at type $v$ by $P_v$, each differential block is a morphism
coefficient times $I_s$, since all of its scalar coefficient matrices
are scalar. The $s$ coordinate lines are therefore invariant and carry
identical typed differentials. Conjugating the usual matrix units back
gives the displayed formula, their multiplication identities, and the
rank-one projections. Their images supply the claimed direct sum. The
same conclusion for extra attachments follows only if the required
invariance or transported block form is checked for those maps as well.
\end{proof}

A sufficient case is that the graph of all nonzero typed blocks is a
tree and each block has a one-dimensional coefficient span whose nonzero
coefficient matrix is invertible. Choosing those matrices as the forest
transports turns every block coefficient matrix into $I_s$. This is an
algebraic family; its prevalence among actual knot-scan components remains
an empirical and structural question. The current implementation computes
and canonically reconstructs the full commutant. It does not yet bypass
the bounded candidate policy with these direct projectors.

Cycles explain a limitation. A non-tree invertible block becomes a
holonomy constraint on $X$, and all noninvertible coefficients impose
their own centralizer conditions. If these matrices have only the scalar
common centralizer, there are no nontrivial projections in that algebra.
Thus a connected invertible forest does not guarantee a useful split.
At the other extreme, low-rank coefficient spans need not display an
invertible basis element: the two singular matrices
$\operatorname{diag}(1,0)$ and $\operatorname{diag}(0,1)$ span $I_2$.
The current forest search can therefore miss an invertible linear
combination. Exhausting the span costs $2^\rho$ candidates, so any wider
search needs a stated rank restriction or its own complexity argument.

\subsection{Explicit and compressed relator growth}

The new shared query costs nearly linearly in the total explicit relator
length $\ell$, with the stated logarithmic and dictionary-model factors.
Its speed does not make $\ell$ small. Generator substitution can duplicate
long words across many relators; a sequence of individually valid moves
can therefore create an expensive explicit presentation before the next
shortening query is reached.

There is a modest unconditional observation for a restricted phase: if
every accepted relator replacement strictly decreases total explicit
length, and no intervening operation increases it, the number of such
replacements is at most the starting length. This says nothing about
whether that phase recognizes every unknot, and it does not bound the
length at entry to later phases. A useful next theorem would control
both presentation growth and the number of verified moves across the
whole producer, including generator eliminations and Whitehead steps.

For straight-line representations, track grammar size, logarithmic
represented length, the number of rewritten nonterminals, and the cost
of checking each move. A shared fully compressed cyclic index would
need to handle donor exclusion and occurrence deletion without expanding
the rotations. An incremental explicit index is another concrete target,
but reusing an index after mutation requires exact removal of obsolete
occurrences. Each proposal should be challenged by presentations where
short descriptions expand rapidly or where most indexed occurrences are
invalidated by a single certified move.

\subsection{Proposed research questions}

The following questions are prioritized by how directly they remove a
known implementation or proof obstruction. They can be pursued separately,
but every claimed global application must satisfy the complete accounting
contract above.

\begin{enumerate}[label=Q\arabic*.,leftmargin=12mm]
\item \textbf{Weighted component extraction.} Can the native interval
kernel return selected component coordinates and grouped topological
weights in polynomial encoded time, including one-sided components and
exponentially large multiplicities?
\item \textbf{Marked cutting.} What is the smallest attachment record that
determines the cut complement and its parallelity bundles relative to all
peripheral markings? Can its expansion be independently compared with
ordinary cutting on every small test input?
\item \textbf{Exterior correspondence.} Can each diagram simplification and
exterior construction emit a short composable certificate preserving the
peripheral identification needed by the final compressing-disc argument?
\item \textbf{Hierarchy repairs.} Is there an amortized potential that bounds
the cumulative cost of discarded and rebuilt hierarchy suffixes by
$\exp(O(\log^2 n))$, or an explicit family refuting a proposed potential?
\item \textbf{Interface compression.} Which exact state equivalence can
reduce the effective interface description below the width of the raw
diagram decomposition while preserving every future attachment operation?
\item \textbf{Natural coefficient ranks.} Are there topologically defined
knot families with uniformly small $\rho_{uv}$ but growing ambient
coefficient support, and can the rank advantage persist after including
all scan and reconstruction costs?
\item \textbf{Constructive scalar summands.} Can the transported-projector
family of Proposition~\ref{prop:scalar-transport-splitting} be recognized
and split directly under the scanner's full attachment contract, avoiding
construction of its entire commutant basis?
\item \textbf{Invertible combinations.} For the coefficient spaces actually
arising in the scanner, can an invertible linear combination be found with
a proved cost smaller than enumerating all $2^\rho$ combinations?
\item \textbf{Canonical output cost.} Can candidate selection be reformulated
in reduced root coordinates with an independently verified equivalence to
the existing policy, so that a large full-basis reconstruction is avoided?
\item \textbf{Shared compressed word search.} Can capped cyclic windows and
distinct-slot donor constraints be indexed in time polynomial in grammar
size and logarithmic represented length, without expanding the rotations?
\item \textbf{Presentation growth.} For which nontrivial diagram families
are both intermediate presentation size and the full verified move count
quasi-polynomially bounded, including the substitutions that increase
length between shortening phases?
\item \textbf{Incremental exact indices.} Can relator replacement and deletion
update the shared index faster than rebuilding it while deleting every
obsolete occurrence and retaining a compact independently checkable match?
\end{enumerate}

\subsection{Milestones that would justify stronger conclusions}

The next report should distinguish the following outcomes.
\begin{enumerate}
\item A completed weighted-orbit and marked-cutting implementation, with
independent reconstruction tests and a proved bound on its encoded output.
\item A verified connection from the ordinary knot diagram to the
triangulation and peripheral data consumed by the geometric kernels.
\item A structural theorem for a nontrivial diagram family controlling
every parameter in Proposition~\ref{prop:global-accounting}, rather than
only the frontier width or a successful witness size.
\item A constructive implementation of Proposition
\ref{prop:scalar-transport-splitting}, checked against the maintained
attachment contract and measured on actual differential components.
\item A complete presentation-growth and move-count bound for a specified
group-simplification family, or an exact compressed replacement for the
steps that defeat such a bound.
\item A hierarchy theorem paying simultaneously for depth, branching,
repairs and representation changes, with every local oracle instantiated.
\end{enumerate}

Any one of these would close a clearly identified gap. A general
$n^{O(\log n)}$ conclusion requires enough of them, or an alternative
complete route, to meet a uniform total-work bound on every input.
